\documentclass[12pt]{article}
\usepackage[spanish,es-tabla]{babel}
%\usepackage{natbib}
\usepackage{url}
\usepackage[utf8]{inputenc}
\usepackage{amsmath}
\usepackage{amssymb} %Para que no webee con que le falta el contador al enumerate
\usepackage{graphicx}
\usepackage{subfigure}
\graphicspath{{images/}}
\usepackage{parskip}
\usepackage{fancyhdr}
\usepackage{vmargin}
\usepackage{gensymb}
\usepackage{adjustbox}
\usepackage{enumerate} %Para que no webee con el enumerate
\usepackage{multirow}%Para usar filas multiples
\usepackage{cite} %Para contraer citas
%\usepackage{hyperref}

\usepackage{listings} %Para incluir matlab -no colocar ñ o tildes en codigo matlab
\usepackage{color} %red, green, blue, yellow, cyan, magenta, black, white
\definecolor{mygreen}{RGB}{28,172,0} % color values Red, Green, Blue
\definecolor{mylilas}{RGB}{170,55,241}


\setmarginsrb{3 cm}{1 cm}{3 cm}{2.5 cm}{1 cm}{1.5 cm}{1 cm}{1.5 cm}

\LARGE{\title{Entrega ...: ...}}								% Title

\author{Álvaro Ignacio Pinto Lohmann}								% Author
\date{\today}											% Date

\makeatletter
\let\thetitle\@title
\let\theauthor\@author
\let\thedate\@date
\makeatother

\pagestyle{fancy}
%\fancyhf{}
%\rhead{\textit{}}
%\lhead{\textit{}}
%\cfoot{\thepage}

\begin{document}

\lstset{language=Matlab,%
    %basicstyle=\color{red},
    breaklines=true,%
    morekeywords={matlab2tikz},
    keywordstyle=\color{blue},%
    morekeywords=[2]{1}, keywordstyle=[2]{\color{black}},
    identifierstyle=\color{black},%
    stringstyle=\color{mylilas},
    commentstyle=\color{mygreen},%
    showstringspaces=false,%without this there will be a symbol in the places where there is a space
    numbers=left,%
    numberstyle={\tiny \color{black}},% size of the numbers
    numbersep=9pt, % this defines how far the numbers are from the text
    emph=[1]{for,end,break},emphstyle=[1]\color{red}, %some words to emphasise
    %emph=[2]{word1,word2}, emphstyle=[2]{style},    
}

%%%%%%%%%%%%%%%%%%%%%%%%%%%%%%%%%%%%%%%%%%%%%%%%%%%%%%%%%%%%%%%%%%%%%%%%%%%%%%%%%%%%%%%%%

\begin{titlepage}
	\centering
    \includegraphics[scale = 0.15]{blank} \hspace{-1 cm} \includegraphics[scale=0.7]{logo_nova}
    
    \vspace{1cm}
	\rule{\linewidth}{0.2 mm} \\[0.4 cm]
	{ \huge \bfseries \thetitle}\\
	\rule{\linewidth}{0.2 mm} \\[1.5 cm]

\vspace{-2cm}
\begin{center}
\end{center}
\begin{center}
\Large{\textbf{ME5701-1 Proyectos en materiales y manufactura}}\\
\end{center}



\vspace{2cm}
Cliente: Sr. Leonel Núñez\\
Equipo de trabajo:\\
Cristobal Bernal\\
Ricardo Fredes\\
Álvaro Pinto\\
Sebastián Venegas\\


\vspace{1 cm}

\textbf{\large{Fecha de entrega:}}
\begin{center}
.... del 2018
\end{center}

\vspace{0.5cm}
NovaManufacturing -Servicios de Ingeniera prestados por empresas N.C.P.\\
Beauchef 851 of 519 Edificio Poniente, Santiago de Chile\\
%Fono: 

\vspace{1 cm}
\vfill
	
\end{titlepage}


\newpage

\section*{Estado de Revisiones}

\begin{table}[h]
\centering
\resizebox{\textwidth}{!}{%

\begin{tabular}{|c|c|c|c|c|c|}
\hline
\multicolumn{6}{|l|}{Estado de Revisiones}\\
\hline
\multicolumn{2}{|l|}{Orden de servicio} & \multicolumn{4}{l|}{NovaManufacturing}\\
\hline
\multicolumn{2}{|l|}{Proyecto}  & \multicolumn{4}{l|}{....}\\
\multicolumn{2}{|l|}{} & \multicolumn{4}{l|}{....}\\
\hline
\multicolumn{2}{|l|}{Título} & \multicolumn{4}{l|}{Informe de selección de ....}\\
\hline
\multicolumn{2}{|l|}{$N^{o}$ de documento} & \multicolumn{4}{l|}{---}\\
\hline
\multicolumn{2}{|l|}{Revisión} & \multicolumn{4}{l|}{Por NovaManufacturing}\\
\hline
$N^{o}$ & Detalle & Datos & Realizó & Revisó & Aprobó\\
\hline
\multirow{4}{*}{A} & \multirow{3}{*}{Entrega de selección de materiales, cubicación y costeo} & Nombre & C.B.-S.V.-A.P.-R.F. & S.V.-A.P. & R.F.\\
\cline{3-6}
& & Fecha & 19-04-2018 & 19-04-2018 & 19-04-2018\\
\cline{3-6}
& & \multirow{12}{*}{ Firma } & \multirow{2}{*}{ } & \multirow{12}{*}{  } & \multirow{12}{*}{  }\\

& &  & &  & \\
& &  & &  & \\
& &  & &  & \\
& &  & &  & \\
& &  & &  & \\
& &  & &  & \\
& &  & &  & \\
& &  & &  & \\
& &  & &  & \\
& &  & &  & \\
& &  & &  & \\
\hline
\end{tabular}%
}

\end{table}


\thispagestyle{empty} %Para que no enumere la pagina

\newpage
\thispagestyle{empty} %Para que no enumere la pagina

\section*{Resumen Ejecutivo}

\qquad En el siguiente informe se presenta ...

\qquad Para la selección de la mayoría de los materiales usados se usó como criterio ....

\qquad Finalmente, se entrega una aproximación de los costos de cada parte de la unidad hidráulica, siendo el costo final de ésta igual a \$.... Un resumen de costos se muestra en la figura siguiente:


%\begin{figure}[h!]
%\centering
%\includegraphics[scale=1]{costos}
%\caption{Resumen de costos de  ...}
%\end{figure}


\newpage
\thispagestyle{empty}
\clearpage
\tableofcontents
\listoffigures
\listoftables
\thispagestyle{empty}
\clearpage
\setcounter{page}{1}

%\section{Utiles}

%\begin{figure}
%\centering
%\includegraphics[]{}
%\label{}\caption{}
%\end{figure}


%\begin{align}
%\frac{}{}   \cdtot \text{kmmlñnxñna.kjqPL}
%\end{align}

\section{Reporte de Avance}

\qquad El proyecto fue entregado el día ... junto a datos pertinentes al mismo como ....

\qquad Desde entonces, NovaManufacturing, soluciones ingenieriles, trabajó en un plan de trabajo a seguir para llevar a cabo el proyecto. Este plan de trabajo se detalla en la carta Gantt y curva S entregadas al cliente el día ... del presente año. Desde entonces se ha trabajado en la selección de materiales, cubicación y cotización para cada pieza que componen las unidades estructurales de ...

\qquad Luego; se entrega un informe el dia .... en donde se detallan las tareas realizadas hasta la fecha.

\qquad A continuacion, en la figura \ref{S_curve}, se muestra el avance en la curva S entregada en el informe actual.

%\begin{figure}[h!]
%\centering
%\includegraphics[scale=1]{curvaS_final}
%\caption{Curva S para fecha ....\label{curvaS}}
%\end{figure}


\newpage

\section{Introducción}

\qquad En el presente informe se da cuenta del estado de avance del proyecto ``....''

\newpage
\section{Revisión de antecedentes pertinentes}

\qquad El objetivo del proyecto dentro del cual se enmarca este informe de avance es estimar el coste de fabricación de ...



\qquad Y se compone de .... subconjuntos principales de componentes:
\begin{itemize}
\item Subconjunto 1: ...
\item Subconjunto 2: ...
\item ...
\end{itemize}





\newpage
\section{Objetivos y alcances}
\subsection{Objetivo General}

\qquad Informar al cliente sobre los materiales seleccionados en la fabricación de..., asi como también sobre la cubicación y costeo de los componentes del prototipo a fabricar.
\subsection{Objetivos Específicos}
\begin{itemize}
\item Analizar y entender los planos para fabricación de ...
\item Analizar los subsistemas que componen ... pieza a pieza.
\item Estudiar los materiales ....
\item Estimar el peso de cada pieza que compone ..., tomando como base los materiales seleccionados y las dimensiones de las piezas.
\item Estimar el costo de fabricación de ...., cotizando el precio por peso de los materiales en el mercado. 
\end{itemize}
\subsection{Alcances}
\qquad ...
\newpage
\section{Metodología}

\qquad El informe presenta la estructura básica para poder ordenar el trabajo en tres subsecciones: selección y justificación de materiales, cubicación y cotización de cada componente. 

\subsection{Selección y justificación de materiales}

\qquad ...

\subsection{Cubicación..}

\qquad ...


\subsection{Cotización}

\qquad ...



\newpage
\section{Resultados}
\subsection{Selección de materiales}

\qquad ... 

\newpage
\subsection{Justificación de material}

\subsubsection{Subconjunto 1}

\qquad ...

\subsubsection{Subconjunto 2}

\qquad ...

\newpage
\subsection{Cubicación}

\qquad ...


\newpage
\subsection{Cotización}

\qquad ...


\newpage
\section{Conclusiones}

\qquad ...


\newpage
\section{Referencias}
\renewcommand{\refname}{}
\begin{thebibliography}{X}

\bibitem{ref1} \textit{...}, ...


\end{thebibliography}

\newpage
\section{Anexos}



\end{document}
